\documentstyle[%


righttag%


,11pt%


,amssymb%


,russian%               remove in Eng.


,izvru%                 Set izven% in Eng.


]{amsart}





% TTS: attn to lines 13 & 14





\newcommand{\I}{\mathrm I}


\newcommand{\II}{\mathrm{II}}


\newcommand{\p}{\mathrm P} % TTS! repl P. for D (see the context, sect. 3)


\newcommand{\C}{\mathrm C} % TTS! Conjugate? (after (51)) or to your taste


\newcommand{\im}{\operatorname{im}}


\newcommand{\pr}{\operatorname{pr}}


\newcommand{\dist}{\operatorname{dist}}


\newcommand{\pen}{\operatorname{pen}}


\newcommand{\tts}[1]{}


\renewcommand{\baselinestretch}{0.97}





\theoremstyle{plain}


\theorembodyfont{\it}


\newtheorem{lemnonum}{\hskip\parindent Лемма}


\renewcommand{\thelemnonum}{}





%\hoffset=-15mm%                  For Eng.


\setcounter{page}{21}


\god{2001}


\nomer{12~(475)} %                        Remove in Eng.


%\nomer{Vol.\,45, No.\,12}%               Set in Eng.


%\str{\pageref{begin}--\pageref{end}}%  Set in Eng.


\udk{519.854}





\author{\it А.И.\,Голиков, Ю.Г.\,Евтушенко}


% NB:   ^^ remove in Eng.


\title{Применение теорем об альтернативах


к нахождению нормальных решений


линейных систем}


\thanks{Работа выполнена при финансовой поддержке Российского фонда


фундаментальных исследований, код проекта \No\,01-01-00804, и по


программе государственной поддержки ведущих научных школ, код


проекта \No\,00-15-96080.}





\date{}


%\pagestyle{myheadings}%         Set in Eng.


\pagestyle{plain} %              Remove in Eng.


\begin{document}


%\thispagestyle{plain}%          Set in Eng.


\maketitle


\label{begin}








\section{Введение}





Численным методам решения систем уравнений и неравенств


посвящена обширная литература. В приложениях часто встречается


случай, когда рассматриваемые системы не имеют решений. В


[1]--[4] такие задачи, названные несобственными, подробно


исследованы и предложены специальные приемы коррекции для


устранения несовместности. Одновременное изучение совместных и


несовместных систем содержится в многочисленных книгах, статьях,


посвященных теоремам об альтернативах (напр., [5]--[10]). В этих


работах каждой системе линейных уравнений и неравенств ставится


в соответствие альтернативная система такая, что обе системы не


могут иметь решения одновременно и всегда разрешима только одна


из них. Априори не известно, имеет ли данная система решениe.


Поэтому следует, во-первых, выяснить, разрешима ли заданная


система и, во-вторых, найти ее решение, если она разрешима.





В данной работе предлагается подход к решению систем линейных


уравнений и неравенств на основе конструктивного доказательства


теорем об альтернативах. Обоснование предлагаемого подхода


проводится на основе теории двойственности нелинейного


программирования. Вводятся системы, сопряженные и альтернативные


к исходной. Размерность переменных в этих системах равна числу


уравнений и неравенств в исходной задаче. Решение сводится к


безусловной минимизации норм невязки исходной или альтернативной


к ней системы.





Если заданы исходная система и альтернативная к ней, то для


нахождения решения разрешимой системы достаточно


проминимизировать невязку какой-либо одной из этих двух систем.


Такая задача дает более богатую информацию, чем просто решение


исходной системы. Если для выбранной системы в результате


безусловной минимизации получена нулевая невязка, то решение


этой задачи одновременно является и решением выбранной системы.


При этом можно утверждать, что вторая система не имеет решения.


Если минимальная невязка выбранной системы не равна нулю, то эта


система не имеет решения, а вторая система разрешима и по


результатам проведенной минимизации по простым формулам


вычисляется решение этой системы с минимальной нормой


(нормальное решение).





Аналогично подходу, используемому в [11], формулируется задача,


которая вместе с задачей безусловной минимизации нормы невязок


составляют пару взаимно двойственных задач. Аппарат теории


двойственности позволяет дать простое конструктивное


доказательство теоремы об альтернативах.





В качестве примера предлагаемого подхода рассматриваются


нетрадиционные двупараметрические условия оптимальности для


задач линейного программирования [12]. Формулируется имеющая


меньшую размерность несовместная система, которая является


альтернативной к системе необходимых и достаточных условий


оптимальности. В результате однократной безусловной минимизации


гладкой кусочно-квадратичной функции с числом переменных на


единицу больше числа переменных прямой задачи вычисляются по


простым формулам нормальное решение прямой задачи и оптимальные невязки


с минимальной нормой двойственной задачи линейного


программирования.





Подчеркнем, что, благодаря различию размерностей переменных


исходной и альтернативной систем, для численного решения


заданной системы может оказаться целесообразным переход от


исходной совместной системы к минимизации невязки альтернативной


несовместной системы. Такая редукция может привести к задаче


безусловной минимизации меньшей размерности и дает возможность


определить нормальное решение исходной системы. Таким образом, в


отличие от устранения несовместности в несобственных задачах


[1]--[4] для нахождения нормального решения совместной системы


может быть полезен переход к ``решению'' несовместной


альтернативной системы.





\section{Основные теоремы}





Пусть $m\times n$-матрица задана в виде


$$


A=


\begin{bmatrix}


A_{11} & A_{12} \\ A_{21} & A_{22}


\end{bmatrix}


,


$$


где прямоугольные матрицы $A_{11}$, $A_{12}$, $A_{21}$, $A_{22}$ имеют 


соответственно размерности


$m_1\times n_1$, $m_1\times n_2$, $m_2\times n_1$, $m_2\times n_2$. Пусть 


векторы $x\in R^n$, $z,u,b\in R^m$ имеют


разбиения $x^\top=[x^\top_1,x^\top_2]$, $z^\top=[z^\top_1,z^\top_2]$, 


$u^\top=[u^\top_1,u^\top_2]$, $b^\top=[b^\top_1,b^\top_2]$, где $x_1\in 


R^{n_1}$, $x_2\in R^{n_2}$, $n=n_1+n_2$, $z_1,u_1,b_1\in R^{m_1}$, 


$z_2,u_2,b_2\in R^{m_2}$, $m=m_1+m_2$.


Введем вспомогательные множества $\Pi_x=\{[x_1,x_2]: x_1\in R^{n_1}_+,\ 


x_2\in R^{n_2}\}$, $\Pi_z=\{[z_1,z_2]:z_1\in R^{m_1}_+,\ z_2\in 


R^{m_2}\}$, $\Pi_u=\{[u_1,u_2]:u_1\in R^{m_1}_+,\ u_2\in R^{m_2}\}$.





Рассмотрим систему линейных равенств и неравенств


\begin{equation}


A_{11}x_1+A_{12}x_2\ge b_1,\quad


A_{21}x_1+A_{22}x_2=b_2,\quad x_1\ge 0_{n_1}.\tag{$\I$}


\end{equation}


Определим {\it сопряженную систему} к $(\I)$


\begin{equation}


A^\top_{11}z_1+A^\top_{21}z_2\le 0_{n_1},\quad


A^\top_{12}z_1+A^\top_{22}z_2=0_{n_2},\quad z_1\ge 0_{m_1}.


\tag{$\I '$}


\end{equation}


Введем {\it альтернативную} к $(\I)$ систему


\begin{equation}


A^\top_{11}u_1+A^\top_{21}u_2\le 0_{n_1},\quad


A^\top_{12}u_1+A^\top_{22}u_2=0_{n_2},\quad


b^\top_1u_1+b^\top_2u_2=\rho,\quad u_1\ge 0_{m_1}.\tag{$\II$}


\end{equation}


Здесь $\rho>0$ --- произвольное фиксированное положительное число. 


Отметим, что


требование положительности $\rho$ автоматически предполагает выполнение 


условия


$\|b\|\ne0$.





Введем вектор $w\in R^{n+1}$, представимый в виде 


$w^\top=[w^\top_1,w^\top_2,w_3]$, где $w_1\in R^{n_1}$, $w_2\in R^{n_2}$, 


$w_3\in R^1$,


и вспомогательное множество $\Pi_w=\{[w_1,w_2,w_3]:w_1\in R^{n_1}_+,\ 


w_2\in R^{n_2},\ w_3\in R^1\}$.





Для системы $(\II)$ сопряженная система имеет вид


\begin{equation}


A_{11}w_1+A_{12}w_2-b_1w_3\ge0_{m_1},\quad


A_{21}w_1+A_{22}w_2-b_2w_3=0_{m_2},\quad


w_1\ge0_{n_1}.\tag{$\II '$}


\end{equation}





Множества решений систем $(\I)$, $(\I ')$, $(\II)$ и $(\II ')$


обозначим соответственно через $X$, $Z$, $U$ и $W$.


В отличие от $(\I)$ и $(\II)$ сопряженные системы $(\I ')$ и $(\II ')$ 


всегда имеют решения,


т.\,к. $0_m\in Z$ и $0_{n+1}\in W$.





\begin{lemnonum}


Системы $(\I)$ и $(\II)$ не могут быть одновременно разрешимы.


\end{lemnonum}





\begin{pf}


Предположим обратное, пусть существуют решения $x^*$ и $u^*$


систем $(\I)$, $(\II)$. Подставим $x^*$ в $(\I)$ и скалярно


умножим первое неравенство в $(\I)$ на $u^*_1$, второе равенство --- на 


$u^*_2$. После


сложения результатов и простейших преобразований получим


$$


x^{*\top}_1(A^\top_{11}u^*_1+A^\top_{21}u^*_2)+


x^{*\top}_2(A^\top_{12}u^*_1+A^\top_{22}u^*_2)\ge


b^\top_1u^*_1+b^\top_2u^*_2.


$$


Согласно $(\II)$ левая часть этого неравенства неположительна, а


правая часть строго положительна, т.\,к. 


$A^\top_{11}u^*_1+A^\top_{21}u^*_2\le0_{n_1}$ и 


$b^\top_1u^*_1+b^\top_2u^*_2=\rho>0$. Приходим к


противоречию.


\end{pf}





Ниже в теоремах 2 и 4 будут приведены простые и конструктивные


доказательства того, что всегда имеет решение одна и только одна


из систем: либо $(\I)$, либо $(\II)$. Поэтому эти системы называются


альтернативными.





Альтернативная к альтернативной задаче $(\II)$ сводится к исходной задаче 


$(\I)$.


Действительно, система, альтернативная к $(\II)$, имеет вид


\begin{equation}


A_{11}w_1+A_{12}w_2-b_1w_3\ge0_{m_1},\quad


A_{21}w_1+A_{22}w_2-b_2w_3=0_{m_2},\quad


\rho w_3=\rho',\ \; w_1\ge0_{n_1},


\end{equation}


где $\rho'>0$ --- произвольное положительное число. Поэтому 


$w_3=\rho'/\rho>0$. Сделав в


(1) замену $x_1=w_1/w_3$, $x_2=w_2/w_3$, приходим к исходной системе 


$(\I)$.





Проекцией точки $\ov x$ на непустое замкнутое множество $X$ назовем точку 


$x^*\in X$,


если $\min\limits_{x\in X}\|\ov x-x\|=\|\ov x-x^*\|$.


Обозначим $x^*=\pr(\ov x,X)$, $\dist(\ov x,X)=\|x^*-\ov x\|$.





Через $\pen(x,X)$ обозначим квадратичный штраф в точке $x$ за нарушение 


условия $x\in X$.


Аналогично введем обозначение $\pen(u,U)$. Таким образом, для множеств


$X$, $U$ и любых $x\in \Pi_x$ и $u\in\Pi_u$ имеем


\begin{align*}


\pen(x,X)&=[\|(b_1-A_{11}x_1-A_{12}x_2)_+\|^2+


\|b_2-A_{21}x_1-A_{22}x_2\|^2]^{1/2},\\


%


\pen(u,U)&=[\|(A^\top_{11}u_1+A^\top_{21}u_2)_+\|^2+


\|A^\top_{12}u_1+A^\top_{22}u_2\|^2+


(\rho-b^\top_1u_1-b^\top_2u_2)^2]^{1/2}.


\end{align*}


Здесь и ниже для простоты считаем, что используется евклидова


норма; $a_+$ есть неотрицательная часть вектора $a$, т.\,е. $i$-я


компонента вектора $a_+$ совпадает с $i$-й компонентой вектора $a$,


если она неотрицательна, и равна нулю в противном случае.





Если $x\in\Pi_x$, то $\pen(x,X)=0$ тогда и только тогда, когда $x\in X$. 


Аналогично,


если $u\in\Pi_u$, то $\pen(u,U)=0$ тогда и только тогда, когда $u\in U$.





Для проверки разрешимости той или иной системы и нахождения


соответствующего решения будем использовать методы безусловной


минимизации, примененные к любой из следующих задач:


\begin{align}


I_1&=\min_{x\in\Pi_x}[\pen(x,X)]^2/2,\\


I_2&=\min_{u\in\Pi_u}[\pen(u,U)]^2/2.


\end{align}





Строго говоря, (2) и (3) не являются задачами безусловной


минимизации, т.\,к. в них присутствуют ограничения на знаки


компонент векторов $x_1$ и $u_1$, однако поскольку большинство методов


безусловной минимизации с помощью простейших модификаций могут


учитывать наличие ограничений на знак переменных, то сохраним за


(2) и (3) этот термин.





Введем две задачи квадратичного программирования


\begin{align}


I^d_1&=\max_{z\in Z}\{b^\top z-\|z\|^2/2\},\\


I^d_2&=\max_{w\in W}\{\rho w_3-\|w\|^2/2\}.


\end{align}





Множества $Z$ и $W$ всегда непусты, т.\,к. содержат нулевые


векторы. В отличие от систем $(\I)$, $(\II)$, которые могут быть


разрешимы или неразрешимы, задачи (2)--(5) всегда имеют решения.


При этом задачи (4) и (5) имеют единственные решения. Задачи (2)


и (3) являются двойственными к задачам (4) и (5) соответственно.


Ниже покажем, что задачи (4) и (5) можно назвать взаимно


двойственными к (2) и (3) соответственно.





Формально задачи безусловной минимизации (2) и (3) не имеют


функций Лагранжа и, следовательно, для них нельзя


непосредственно построить двойственные задачи. Тем не менее с


помощью введения дополнительных переменных можно построить


искусственные ограничения и получить эквивалентные задачи


нелинейного программирования, для которых определены


двойственные задачи. Такой прием, насколько известно авторам,


был впервые четко сформулирован и использован в [11].





Введем вектор дополнительных переменных $y\in R^m$, 


$y^\top=[y^\top_1,y^\top_2]$, где $y_1\in R^{m_1}$, $y_2\in R^{m_2}$,


такой, что


$$


y_1=b_1-A_{11}x_1-A_{12}x_2,\quad


y_2=b_2-A_{21}x_1-A_{22}x_2.


$$


Тогда задача (2) заменяется эквивалентной задачей на условный минимум


\begin{equation}


I_1=\min_{[x,y]\in G}f(y),


\end{equation}


где целевая функция и допустимое множество $G$ имеют вид


\begin{gather*}


f(y)=\|(y_1)_+\|^2/2+\|y_2\|^2/2,\\


G=\{[x,y]:A_{11}x_1+A_{12}x_2+y_1=b_1,


\ A_{21}x_1+A_{22}x_2+y_2=b_2,\ x\in\Pi_x\}.


\end{gather*}





В отличие от множества $X$ множество $G$ всегда непусто.





Для задачи квадратичного программирования (6) определим вектор множителей


Лагранжа $z\in\Pi_z$ и выпишем функцию Лагранжа


$$


L(x,y,z)=f(y)+z^\top_1(b_1-A_{11}x_1-A_{12}x_2-y_1)+


z^\top_2(b_2-A_{21}x_1-A_{22}x_2-y_2).


$$


Преобразуем ее к виду


\begin{equation}


L(x,y,z)=f(y)-x^\top_1(A^\top_{11}z_1+A^\top_{21}z_2)-


x^\top_2(A^\top_{12}z_1+A^\top_{22}z_2)+


z^\top_1(b_1-y_1)+z^\top_2(b_2-y_2).


\end{equation}





Введем двойственную функцию


\begin{equation}


F(z)=\min_{x\in\Pi_x}\; \min_{y\in R^n}L(x,y,z)


\end{equation}


и рассмотрим двойственную к (6) задачу $\max\limits_{z\in\Pi_z}F(z)$.





Необходимые и достаточные условия минимума в задаче (8) имеют вид


\begin{align}


L_{x_1}(x,y,z)&=-A^\top_{11}z_1-A^\top_{21}z_2\ge0_{n_1},\ \;


D(x_1)(A^\top_{11}z_1+A^\top_{21}z_2)=0_{n_1},\ \; x_1\ge0_{n_1},\\


%


L_{x_2}(x,y,z)&=-A^\top_{12}z_1-A^\top_{22}z_2=0_{n_2},\\


%


L_{y_1}(x,y,z)&=(y_1)_+-z_1=0_{m_1},\quad


L_{y_2}(x,y,z)=y_2-z_2=0_{m_2}.


\end{align}


Здесь и ниже через $D(z)$ обозначена диагональная матрица, у которой $i$-й


диагональный элемент есть $i$-я компонента вектора $z$.





При $z\in\Pi_z$ из (11) имеем $z=y$. Подставим это соотношение в (7) и 


потребуем


выполнения условия $z\in Z$. Тогда согласно (8)--(10) двойственная


функция $F(z)$ принимает вид $F(z)=b^\top z-\|z\|^2/2$. Таким образом, 


приходим к


задаче (4), которая является двойственной к (6) и,


следовательно, к (2). Итак, задачу безусловной минимизации (2) и


задачу квадратичного программирования (4) можно назвать взаимно


двойственными. Аналогично задачи (3) и (5) являются


двойственными друг к другу.





\begin{thm}


Всякое решение $x^*$ задачи $(2)$ определяет единственное решение


$z^{*\top}=[z^{*\top}_1,z^{*\top}_2]$ задачи $(4)$ по формулам


\begin{equation}


z^*_1=(b_1-A_{11}x^*_1-A_{12}x^*_2)_+,\quad


z^*_2=b_2-A_{21}x^*_1-A_{22}x^*_2,


\end{equation}


и справедливы следующие соотношения$:$


\begin{gather}


\|z^*\|^2=b^\top z^*,\\


z^*\bot Ax^*,\quad z^*\bot(b-z^*),\\


z^*=\pr(b,Z),\quad \|z^*\|=\pen(x^*,X),\quad


\|b-z^*\|=\dist(b,Z),\\


[\pen(x^*,X)]^2+[\dist(b,Z)]^2=\|b\|^2.


\end{gather}


\end{thm}





\begin{pf}


В точке $x^*$ необходимые и достаточные условия минимума


задачи (2) записываются в виде


\begin{gather}


-A^\top_{11}(b_1-A_{11}x^*_1-A_{12}x^*_2)_+-


A^\top_{21}(b_2-A_{21}x^*_1-A_{22}x^*_2)\ge0_{n_1},\notag\\


D(x^*_1)[A^\top_{11}(b_1-A_{11}x^*_1-A_{12}x^*_2)_++


A^\top_{21}(b_2-A_{21}x^*_1-A_{22}x^*_2)]=0_{n_1},


\ \; x^*_1\ge0_{n_1},\\


A^\top_{12}(b_1-A_{11}x^*_1-A_{12}x^*_2)_++


A^\top_{22}(b_2-A_{21}x^*_1-A_{22}x^*_2)=0_{n_2}.\notag


\end{gather}





В формулах (17) введем обозначения


\begin{equation}


z^*_1=(b_1-A_{11}x^*_1-A_{12}x^*_2)_+,\quad


z^*_2=b_2-A_{21}x^*_1-A_{22}x^*_2


\end{equation}


и покажем, что $z^{*\top}=[z^{*\top}_1,z^{*\top}_2]$ является решением 


задачи (4).


Условия (17) запишем в виде


\begin{gather}


A^\top_{11}z^*_1+A^\top_{21}z^*_2\le0_{n_1},\quad


A^\top_{12}z^*_1+A^\top_{22}z^*_2=0_{n_2},\\


D(x^*_1)(A^\top_{11}z^*_1+A^\top_{21}z^*_2)=0_{n_1},\quad


x^*_1\ge0_{n_1}.


\end{gather}





Из (18) и (19) следует, что $z^*\in Z$. Умножим в (18) первое соотношение 


скалярно


на $z^*_1$, второе --- на $z^*_2$, получим


\begin{align*}


\|z^*_1\|^2&=z^{*\top}


_1(b_1-A_{11}x^*_1-A_{12}x^*_2)_+=z^{*\top}_1(b_1-A_{11}x^*_1-


A_{12}x^*_2)=


b^\top_1 z^*_1-x^{*\top}_1A^\top_{11}z^*_1-x^{*\top}_2


A^\top_{12}z^*_1,\\


%


\|z^*_2\|^2&=z^{*\top}_2(b_2-A_{21}x^*_1-A_{22}x^*_2)=


b^\top_2 z^*_2-x^{*\top}_1A^\top_{21}z^*_2-


x^{*\top}_2A^\top_{22}z^*_2.


\end{align*}


Складывая эти соотношения, находим


\begin{multline}


\|z^*\|^2=\|z^*_1\|^2+\|z^*_2\|^2=b^\top_1 z^*_1+


b^\top_2z^*_2-x^{*\top}_1(A^\top_{11}z^*_1+A^\top_{21}z^*_2)-


x^{*\top}_2(A^\top_{12}z^*_1+A^\top_{22}z^*_2)=\\


=(b-Ax^*)^\top z^*=b^\top z^*.


\end{multline}


В (21) учтено, что согласно (19) и (20)


$x^{*\top}_1(A^\top_{11}z^*_1+A^\top_{21}z^*_2) 


+x^{*\top}_2(A^\top_{12}z^*_1+A^\top_{22}z^*_2)= x^{*\top}A^\top z^*=0$.


Таким образом, доказано равенство (13) и показано, что $z^{*\top}Ax^*=0$,


т.\,е. векторы $z^*$ и $Ax^*$ ортогональны. Условие (13) можно переписать


в виде $z^{*\top}(z^*-b)=0$, откуда следует второе из утверждений (14).





Запишем функцию Лагранжа для задачи (4)


\begin{equation}


L(z,x)=b^\top z-\|z\|^2/2-x^\top_1(A^\top_{11}z_1+


A^\top_{21}z_2)-x^\top_2(A^\top_{12}z_1+A^\top_{22}z_2)


\end{equation}


и приведем условия Куна--Таккера


\begin{gather}


L_{z_1}(z,x)=b_1-z_1-A_{11}x_1-A_{12}x_2\le0_{m_1},\\


%


D(z_1)(b_1-z_1-A_{11}x_1-A_{12}x_2)=0_{m_1},\quad z_1\ge0_{m_1},\\


%


L_{z_2}(z,x)=b_2-z_2-A_{21}x_1-A_{22}x_2=0_{m_2},\\


%


L_{x_1}(z,x)=-(A^\top_{11}z_1+A^\top_{21}z_2)\ge0_{n_1},


\ \; x_1\ge0_{m_1},\ \; D(x_1)(A^\top_{11}z_1+


A^\top_{21}z_2)=0_{n_1},\\


%


L_{x_2}(z,x)=-(A^\top_{12}z_1+A^\top_{22}z_2)=0_{n_2}.


\end{gather}





Сравним необходимые и достаточные условия минимума (18)--(20) для


задачи (2) с необходимыми и достаточными условиями оптимальности


(23)--(27) для задачи квадратичного программирования (4). Если в


условиях Куна--Таккера (23)--(27) в качестве $x$ взять вектор $x^*$,


а в качестве $z$ --- вектор $z^*$, определяемый в (18), то (26),


(27) переходят в (19), (20). Легко видеть, что соотношения (18)


обеспечивают выполнение условий (23)--(25). Таким образом,


седловую точку $[z^*,x^*]$ функции Лагранжа (22) составляют вектор $z^*$ 


---


решение задачи (4) и вектор $x^*$ --- решение задачи (2), причем


между этими векторами имеется связь, задаваемая формулами (12).





Из (2), (4) и (18) следует


$I^d_1=\|z^*\|^2/2=I_1=[\pen(x^*,X)]^2/2$.


Отсюда получаем второе из утверждений (15).





Задачу (4) преобразуем к двум эквивалентным задачам


$$


I^d_1=\max_{z\in Z}[-\|b-z\|^2+\|b\|^2]/2=


\|b\|^2/2-\min_{z\in Z}\|b-z\|^2/2.


$$


Так как $z^*$ есть единственное решение задачи (4), то приходим


к выводу, что $z^*=\pr(b,Z)$ и $\|b-z^*\|=\dist(b,Z)$, т.\,е. доказаны все 


утверждения (15).





Из (13) следует, что вектор невязки $z^*$ ортогонален вектору $b-z^*$.


Поэтому три точки: начало координат в $R^m$, точки $z^*$ и $b$ образуют


прямоугольный треугольник, у которого вектор $b$ определяет


гипотенузу, вектор $z^*$ --- катет длиной, равной $\pen(x^*,X)$, вектор 


$b-z^*$ ---


второй катет длиной $\dist(b,Z)$. Из теоремы Пифагора получим (16).


\end{pf}





Из теоремы 1 следует, что $\|z^*\|\le\|b\|$, $b^\top z^*\ge0$. Вектор 


$z^*$ лежит в полусфере


с центром в начале координат и радиусом, равным $\|b\|$, причем $z^*$ 


лежит в той


половине сферы, где вектор $z^*$ образует острый угол с вектором $b$.





В теореме 1 утверждение (13) является следствием двойственности задач (4) 


и


(2). Действительно, условия (12) позволяют выразить оптимальные переменные


$x^*$ в задаче (2) через оптимальные переменные $z^*$ задачи (4). Тогда из 


условия


двойственности $I_1=I^d_1$ получаем $\|z^*\|^2=b^\top z^*$.





Теперь можно сформулировать и доказать следующую теорему об альтернативах.





\begin{thm}


Пусть $x^{*\top}=[x^{*\top}_1,x^{*\top}_2]$ --- произвольное решение 


задачи $(2)$, $z^{*\top}=[z^{*\top}_1,z^{*\top}_2]$ ---


единственное решение задачи $(4)$.


\begin{itemize}


\item[1.] Если $\|z^*\|=0$, то $I_1=I^d_1=0$, $\dist(b,Z)=\|b\|$ и 


система $(\I)$ разрешима,


одним из ее решений является $x^*;$ система $(\II)$ неразрешима.


\item[2.] Если $\|z^*\|\ne0$, то $I_1=I^d_1>0$, $\dist(b,Z)<\|b\|$ и 


система $(\I)$ неразрешима$;$


система $(\II)$ разрешима и вектор $u^*=\rho z^*/\|z^*\|^2$ --- ее 


нормальное решение.


\end{itemize}


\end{thm}





\begin{pf}


1. Если $\|z^*\|=0$, то из (18) следует, что система $(\I)$ разрешима и 


$x^*$ --- одно


из решений системы $(\I)$, при этом $I_1=I^d_1=0$ и (16) переходит в 


выражение


$\dist(b,Z)=\|b\|$. Согласно лемме система $(\II)$ не имеет решения.





2. Пусть $\|z^*\|\ne0$. Тогда из (12) следует, что система $(I)$ 


неразрешима и


$I_1=I^d_1>0$. Так как $I_1=\pen(x^*,X)>0$, то из (16) имеем 


$\dist(b,Z)<\|b\|$.





В процессе доказательства предыдущей теоремы было показано, что вектор


$z^*$ удовлетворяет сопряженной однородной системе $(\I ')$. Поэтому 


вектор $u^*=\rho z^*/\|z^*\|^2$


также удовлетворяет системе $(\I ')$. По сравнению с $(\I ')$ в $(\II)$ 


добавлено


одно неоднородное уравнение. Подставляя $u^*$ в это уравнение, 


убеждаемся с


учетом (13), что вектор $u^*\in U$. Покажем, что $u^*$ --- нормальное 


решение системы


$(\II)$, т.\,е. является решением задачи


\begin{equation}


\min_{u\in U}\|u\|^2/2.


\end{equation}


Двойственной к (28) является задача


\begin{equation}


\max_{\Hat x_1\in R^{n_1}_+}\; \max_{\Hat x_2\in R^{n_2}}


\; \max_{\Hat x_3\in R^1}[\rho\wh x_3-\|(b_1\wh x_3-A_{11}\wh x_1-


A_{12}\wh x_2)_+\|^2/2-\|b_2\wh x_3-A_{21}\wh x_1-


A_{22}\wh x_2\|^2/2].


\end{equation}


Запишем функцию Лагранжа для задачи (28)


$$


L(u,\wh x)=\|u\|^2/2+\wh x^\top_1(A^\top_{11}u_1+A^\top_{21}u_2)+


\wh x^\top_2(A^\top_{12}u_1+A^\top_{22}u_2)+


\wh x_3(\rho-b^\top_1u_1-b^\top_2u_2)


$$


и приведем условия Куна--Таккера, вычисленные в седловой точке $[\wt 


u^*,\wh x^*]$, где $\wt u^{*\top}=[\wt u^{*\top}_1,\wt u^{*\top}_2]$


--- решение задачи (28) и $\wh x^{*\top}=[\wh x^{*\top}_1,\wh 


x^{*\top}_2]$ --- решение задачи (29),


\begin{gather}


\wt u^*_1+A_{11}\wh x^*_1+A_{12}\wh x^*_2-b_1\wh x^*_3\ge0_{m_1},\\


%


D(\wt u^*_1)(\wt u^*_1+A_{11}\wh x^*_1+A_{12}\wh x^*_2-


b_1\wh x^*_3)=0_{m_1},\ \; \wt u^*_1\ge0_{m_1},\\


%


\wt u^*_2+A_{21}\wh x^*_1+A_{22}\wh x^*_2-


b_2\wh x^*_3=0_{m_2},\\


%


A^\top_{11}\wt u^*_1+A^\top_{21}\wt u^*_2\le0_{n_1},\quad


D(\wh x^*_1)(A^\top_{11}\wt u^*_1+A^\top_{21}\wt u^*_2)=


0_{n_1},\quad \wh x^*_1\ge0_{n_1},\\


%


A^\top_{12}\wt u^*_1+A^\top_{22}\wt u^*_2=0_{n_2},\\


%


\rho-b^\top_1\wt u^*_1-b^\top_2\wt u^*_2=0.


\end{gather}





Из (30)--(32) получаем, что решение $\wt u^*$ задачи (28) связано с 


решением $\wh x^*$ задачи


(29) соотношениями


$$


\wt u^*_1=(b_1\wh x^*_3-A_{11}\wh x^*_1-A_{12}\wh x^*_2)_+,\quad


\wt u^*_2=b_2\wh x^*_3-A_{21}\wh x^*_1-A_{22}\wh x^*_2.


$$


Используя эти соотношения, из равенств оптимальных значений целевых 


функций


прямой задачи (28) и двойственной (29) имеем $\|\wt u^*\|^2=\rho\wh 


x^*_3$. Отсюда следует $\wh x^*_3>0$.





Сделав в (30)--(34) замену $\wt u^*=\wh x^*_3z^*$, $\wh x^*_1=\wh 


x^*_3x^*_1$, $\wh x^*_2=\wh x^*_3x^*_2$


и сократив эти выражения на положительную величину $\wh x^*_3$, приходим к 


условиям


Куна--Таккера (23)--(27) для задачи (2), выполненным в точке $[z^*,x^*]$. 


Подставляя


$\wt u^*=\wh x^*_3z^*$ в (35), с учетом (13) получаем


$$


\rho/\wh x^*_3-b^\top z^*=\rho/\wh x^*_3-\|z^*\|^2.


$$


Отсюда при $\wh x^*_3=\rho/\|z^*\|^2$ получаем, что нормальное решение 


системы $(\II)$ выражается


через решение задачи (2) по формуле $\wt u^*=\rho z^*/\|z^*\|^2$.


\end{pf}





Анализ задач (3) и (5) аналогичен приведенному для задач (2),


(4), однако есть некоторые отличительные стороны. Поэтому о


свойствах задач (3), (5) приведем две теоремы, аналогичные


теоремам 1, 2. Введем вектор $r\in R^{n+1}$, заданный в виде 


$r^\top=[0^\top_n,\rho]$, и матрицу $\wh A=[-A\mid b]$.





\begin{thm}


Пусть $u^{*\top}=[u^{*\top}_1,u^{*\top}_2]$ --- произвольное решение 


задачи $(3)$. Тогда


решение $w^{*\top}=[w^{*\top}_1,w^{*\top}_2,w^*_3]$ задачи $(5)$ 


выражается через решение $w^*$ задачи $(3)$


по формулам


\begin{equation}


w^*_1=(A^\top_{11}u^*_1+A^\top_{21}u^*_2)_+,\quad


w^*_2=A^\top_{12}u^*_1+A^\top_{22}u^*_2,\quad


w^*_3=\rho-b^*_1u^*_1-b^\top_2u^*_2,


\end{equation}


причем вектор $w^*$ обладает следующими свойствами$:$


\begin{gather}


\|w^*\|^2=\rho w^*_3; \\


%


w^*\bot \wh A^\top u^*,\quad w^*\bot(r-w^*),\\


%


w^*=\pr(r,W),\ \; \|w^*\|=\pen(u^*,U),


\ \; \|r-w^*\|=\dist(r,W), \\


[\pen(u^*,U)]^2+[\dist(r,W)]^2=\|r\|^2,\\


%


\|w^*\|\le\rho,\ \; 0\le w^*_3\le\rho,


\ \; \|w^*_1\|^2+\|w^*_2\|^2\le\rho^2/4.


\end{gather}


\end{thm}





\begin{pf}


Задача строго выпуклого квадратичного программирования (5) является 


задачей


нахождения проекции вектора $[0_n,\rho]$ на непустое множество $W$, 


заданное


системой линейных равенств и неравенств $(\II ')$. Эта задача всегда имеет


единственное решение, для нее существует вектор множителей Лагранжа 


$u\in R^n$,


а функция Лагранжа имеет вид


$$


L(w,u)=\rho w_3-\|w\|^2/2-u^\top_1(b_1w_3-A_{11}w_1-A_{12}w_2)-


u^\top_2(b_2w_3-A_{21}w_1-A_{22}w_2).


$$


Запишем необходимые и достаточные условия оптимальности (условия


Куна--Таккера) для задачи (5), вычисленные в седловой точке $[w^*,u^*]$,


\begin{gather}


-w^*_1+A^\top_{11}u^*_1+A^\top_{21}u^*_2\le0_{n_1},\quad


D(w^*_1)(-w^*_1+A^\top_{11}u^*_1+A^\top_{21}u^*_2)=0_{n_1},


\ \; w^*_1\ge0_{n_1},\\


%


-w^*_2+A^\top_{12}u^*_1+A^\top_{22}u^*_2=0_{n_2},\\


%


\rho-w^*_3-b^\top_1 u^*_1-b^\top_2u^*_2=0,\\


%


A_{11}w^*_1+A_{12}w^*_2-b_1w^*_3\ge0_{m_1},\quad


D(u^*_1)(A_{11}w^*_1+A_{12}w^*_2-b_1w^*_3)=0_{m_1},


\ \; u^*_1\ge0_{m_1},\\


%


A_{21}w^*_1+A_{22}w^*_2-b_2w^*_3=0_{m_2}.


\end{gather}





Из условий (42)--(44) следуют формулы (36) для определения вектора $w^*$.


Подставляя их в условия (45), (46), приходим к соотношениям


\begin{gather*}


A_{11}(A^\top_{11}u^*_1+A^\top_{21}u^*_2)_++A_{12}


(A^\top_{12}u^*_1+A^\top_{22}u^*_2)-b_1(\rho-b^\top_1u^*_1-


b^\top_2u^*_2)\ge0_{m_1},\\


%


D(u^*_1)[A_{11}(A^\top_{11}u^*_1+A^\top_{21}u^*_2)_++


A_{12}(A^\top_{12}u^*_1+A^\top_{22}u^*_2)-


b_1(\rho-b^\top_1u^*_1-b^\top_2u^*_2)]=0_{m_1},


\ \; u^*_1\ge0_{m_1}, \\


%


A_{21}(A^\top_{11}u^*_1+A^\top_{21}u^*_2)_++A_{22}


(A^\top_{12}u^*_1+A^\top_{22}u^*_2)-


b_2(\rho-b^\top_1u^*_1-b^\top_2u^*_2)=0_{m_2}.


\end{gather*}


Эти соотношения совпадают с необходимыми и достаточными


условиями оптимальности задачи (3), вычисленными в точке $u^*$.


Итак, из условий Куна--Таккера и условий оптимальности задачи


(3) следует, что в седловой точке $[w^*,u^*]$ вектор $w^*$ --- решение


задачи (5) и $u^*$ --- решение задачи (3). Эти решения связаны


соотношениями (36).





Из равенства целевых функций прямой (5) и двойственной (3) задач с учетом


соотношений (36) получаем (37).





Доказательство соотношений (38)--(40) аналогично доказательству


соотношений (14)--(16) в теореме 1.





При $\|w^*\|\ne0$ из (37) следует, что $w^*_3>0$. Рассмотрим (37) как 


квадратное


уравнение относительно $w^*_3$


\begin{equation}


w^{*^2}_3-\rho w^*_3+\|w^*_1\|^2+\|w^*_2\|^2=0.


\end{equation}





Из того, что $\|r\|=\rho$ и вектор $w^*$ есть проекция $r$ на $W$, следует


первое неравенство в (41). Вторая оценка в (41) следует из


условия неотрицательности свободного члена квадратного


уравнения $w^*_3(\rho-w^*_3)=\|w^*_1\|^2+\|w^*_2\|^2\ge0$, а третья 


--- из условия неотрицательности


детерминанта квадратного уравнения (47).


\end{pf}





\begin{thm}


Пусть $u^{*\top}=[u^{*\top}_1,u^{*\top}_2]$


--- произвольное решение задачи $(3)$,


$w^{*\top}=[w^{*\top}_1,w^{*\top}_2,w^*_3]$ ---


решение задачи $(5)$, определяемое по формулам $(36)$.


\begin{itemize}


\item[1.] Если $\|w^*\|=0$, то $I_2=I^d_2=0$, система $(\II)$ разрешима, 


одним из ее решений


является $u^*$, система $(\I)$ неразрешима.


\item[2.] Если $\|w^*\|\ne0$, то $w^*_3>0$, $I_1=I^d_1>0$, система 


$(\II)$ неразрешима$;$ система


$(\I)$ разрешима и вектор $x^*_1=w^*_1/w^*_3$, $x^*_2=w^*_2/w^*_3$ --- ее 


решение с минимальной


евклидовой нормой.


\end{itemize}


\end{thm}





\begin{pf}


1. Если $w^*=0_{n+1}$, то из (36) следует, что система $(\II)$ разрешима и 


одним


из ее решений является вектор $u^*$. Согласно лемме система $(\I)$ не 


может быть


разрешима.





2. Пусть $w^*\ne0_{n+1}$. Тогда из (36) следует, что система $(\II)$ 


неразрешима и


$I_2=I^d_2>0$. При этом из условия (37) получаем $w^*_3>0$. Из условий 


(45), (46)


следует, что $x^*_1=w^*_1/w^*_3$, $x^*_2=w^*_2/w^*_3$ является решением 


системы $(\I)$ и имеет минимальную евклидову норму, т.\,е. является решением


задачи


\begin{equation}


\min_{x\in X}\|x\|^2.


\end{equation}





Приведем функцию Лагранжа для этой задачи


$$


L(x,\wt u)=\|x\|^2/2+\wt u^\top_1(b_1-A_{11}x_1-A_{12}x_2)+


\wt u^\top_2(b_2-A_{21}x_1-A_{22}x_2)


$$


и условия Куна--Таккера, вычисленные в седловой точке $[x^*,\wt u^*]$, 


где $x^{*\top}=[x^{*\top}_1,x^{*\top}_2]$


--- решение задачи (48),


$\wt u^{*\top}=[\wt u^{*\top}_1,\wt u^{*\top}_2]$


--- оптимальный вектор множителей


Лагранжа в задаче (48),


\begin{gather*}


x^*_1-A^\top_{11}\wt u^*_1-A^\top_{21}\wt u^*_2\ge0_{n_1},\quad


D(x^*_1)(x^*_1-A^\top_{11}\wt u^*_1-A^\top_{21}\wt u^*_2)=0_{n_1},


\quad x^*_1\ge0_{n_1}, \\


%


x^*_2-A^\top_{12}\wt u^*_1-A_{22}\wt u^*_2=0_{n_2},\\


%


b_1-A_{11}x^*_1-A_{12}x^*_2\le0_{m_1},\quad


D(\wt u^*_1)(b_1-A_{11}x^*_1-A_{12}x^*_2)=0_{m_1},\quad


\wt u^*_1\ge0_{m_1},\\


%


b_2-A_{21}x^*_1-A_{22}x^*_2=0_{m_2}.


\end{gather*}


Эти условия Куна--Таккера для задачи (48) получаются из


(42), (43), (45), (46), если


положить $x^*_1=w^*_1/w^*_3$, $x^*_2=w^*_2/w^*_3$, $\wt 


u^*_1=u^*_1/w^*_3$, $\wt u^*_2=u^*_2/w^*_3$.


\end{pf}





Возможны различные варианты представления альтернативной системы


$(\II)$. Как следует из приведенных теорем, система, альтернативная к


$(\I)$, получается из сопряженной системы $(\I ')$ путем добавления


условия, исключающего возможность тривиального решения.


Например, можно потребовать для решений сопряженной системы $(\I ')$


выполнения условия $b^\top u>0$ (как в лемме Фаркаша [6]) либо условия 


$b^\top u=1$


(как в теореме Гейла [5]), либо наложить нелинейное условие 


$\|u\|^2=\rho$,


где $\rho>0$ --- произвольная фиксированная величина, и т.\,д.





Приведем геометрическую интерпретацию полученных результатов. Пусть 


$b^\bot$


обозначает проекцию вектора $b$ на множество $Z$. Тогда ортогональный к 


нему


вектор будет $b^\parallel=b-b^\bot$, соотношения (15) запишутся в виде


$$


z^*=b^\bot=\pr(b,Z),\quad \|b^\bot\|=\pen(x^*,X),\quad


\|b^\parallel\|=\dist(x^*,X).


$$





В этих обозначениях соотношения (13) и (16) становятся


очевидными: $\|b^\bot\|^2=\|b^\bot\|^2$, 


$\|b^\bot\|^2+\|b^\parallel\|^2=\|b\|^2$.


Аналогично в теореме 3 имеем


$$


w^*=r^\bot=\pr(r,W),\quad \|r^\bot\|=\pen(u^*,U),\quad


r^\parallel=r-r^\bot,\quad


\|r^\parallel\|=\dist(r,W).


$$


Очевидно, векторы $r^\bot$ и $r^\parallel$ лежат внутри шара радиуса 


$\rho$ с центром в начале


координат.





\section{Приложение к линейному программированию}





Применим приведенные выше результаты к задачам линейного


программирования (ЛП). Пусть прямая задача ЛП задана в каноническом виде


\begin{equation}


\min_{x\in X}c^\top x,\quad X=\{x\in R^n:Ax=b,\ x\ge0_n\}.\tag{$\p$}


\end{equation}


Здесь и ниже $A$ --- $m\times n$-матрица ранга $m$, $m<n$, дефект матрицы 


$A$, равный


$n-m$, обозначим через $\nu$, векторы $c,x\in R^n$, $b\in R^m$.





Ниже вместо традиционных необходимых и достаточных условий оптимальности


для задач ЛП воспользуемся условиями из [12]. Для этого введем матрицу


$K$ размера $\nu\times n$. В качестве $K$ можно использовать любую 


матрицу, $\nu$ строк


которой образуют базис нуль-пространства матрицы $A$. Таким образом, 


$\im K^\top$


является ортогональным дополнением к пространству $\im A^\top$. Поэтому


\begin{equation}


\im K^\top=\ker A,\quad A K^\top=0_{m\nu},\quad


R^n=\im A^\top\oplus \im K^\top.


\end{equation}


Здесь через $0_{ij}$ обозначена $i\times j$-матрица с нулевыми элементами.





Если матрицу $A$ представить в блочном виде $A=[B\mid N]$, где $B$ 


невырождена,


то матрицу $K$ можно записать в виде $K=[-N^\top(B^{-1})^\top\mid I_\nu]$. 


Если с помощью


преобразований Гаусса--Жордана матрицу $A$ привести к виду $A=[I_m\mid 


N]$,


тогда матрица $K$ представима в виде $K=[-N^\top\mid I_\nu]$. В ряде 


случаев построение


матрицы $K$ упрощается. Например, пусть задача ЛП имеет ограничения типа


неравенств $N z\ge b$, где $N$ --- $m\times\nu$-матрица, $z\in R^\nu_+$. 


Вводя дополнительные


переменные $\xi\in R^m_+$, вектор $x\in R^n$ представим как объединение 


векторов $z$, $\xi$,


т.\,е. $x^\top=[z^\top,\xi^\top]$. Поэтому допустимое множество запишется в 


том же виде, что и в


задаче $(\p)$, где матрица $A=[N\mid-I_m]$. Тогда $K=[I_\nu\mid N^\top]$.





Определим вектор $d=K c\in R^\nu$ и вектор невязок


\begin{equation}


v=c-A^\top u.


\end{equation}





Введем два аффинных множества


$$


\ov X=\{x\in R^n:Ax=b\},\quad


\ov V=\{v\in R^n:K v=d\}.


$$


Всюду ниже через $\ov x$ и $\ov v$ будем обозначать произвольные


фиксированные $n$-мерные векторы, удовлетворяющие соответственно


условиям $\ov x\in\ov X$, $\ov v\in\ov V$. Заметим, что некоторые 


компоненты векторов $\ov x$,


$\ov v$ могут быть отрицательными. В простейшем варианте можно взять


$\ov v=c$. В силу того, что матрица $A$ имеет ранг $m$ и $n>m$, всегда 


$\ov X\ne\emptyset$ и


$\ov V\ne\emptyset$. Для нахождения вектора $\ov v\in \ov V$ достаточно 


взять произвольный


вектор $\ov u\in R^m$ и положить $\ov v=c-A^\top\ov v$. Принимая во 


внимание (49) и (50),


приходим к выводу, что $\ov v\in\ov V$.





Следуя [12], необходимые и достаточные условия минимума в задаче $(\p)$


запишем в виде системы линейных уравнений с неотрицательными переменными


\begin{equation}


\begin{bmatrix}


A & 0_{mn} \\ 0_{\nu n} & K \\ \ov v^\top & \ov x^\top


\end{bmatrix}


\begin{bmatrix}


x \\ v


\end{bmatrix}


=


\begin{bmatrix}


b \\ d \\ \ov x^\top\ov v


\end{bmatrix}


,\quad x\ge0_n,\quad v\ge 0_n.


\end{equation}





Если задача ЛП $(\p)$ имеет решение, то система (51) совместна и, решив 


ее,


находим решение прямой задачи $(\p)$ и сопряженной задачи ЛП


\begin{equation}


\min_{v\in V}\ov x^\top v,\quad V=\{v\in R^n:K v=d,\ v\ge0_n\}.


\tag{$\C$}


\end{equation}


Система (51) состоит из $n+1$ условий равенств, $2n$ неравенств и


содержит $2n$ переменных. Из-за высокой размерности ее


непосредственное решение может оказаться затруднительным.


Поэтому перейдем к альтернативной системе, которая имеет меньше


переменных, чем система (51), и состоит из $2n$ линейных неравенств


и одного равенства с $n+1$ переменными


\begin{equation}


\begin{bmatrix}


A^\top & 0_{n\nu} & \ov v \\ 0_{nm} & K^\top & \ov x


\end{bmatrix}


\begin{bmatrix}


p \\ q \\ \alpha


\end{bmatrix}


\le0_{2n},\quad b^\top p+d^\top q+\ov x^\top\ov v\alpha=\rho.


\end{equation}


Здесь $\rho>0$ --- произвольная фиксированная константа, $p\in R^m$, 


$q\in R^\nu$, $\alpha\in R^1$.





Так как система (51) совместна, то альтернативная система (52)


несовместна. Это позволяет на основании теорем 3, 4 из


безусловной минимизации невязок системы (52) получить нормальные


решения задач $(\p)$ и $(\C)$. Задача (3) в данном случае записывается в


виде


\begin{equation}


\min_{p\in R^m}\;\min_{q\in R^\nu}\; \min_{\alpha\in R^1}


[\|(A^\top p+\ov v\alpha)_+\|^2+


\|(K^\top q+\ov x\alpha)_+\|^2+


(\rho-b^\top p-d^\top q-\ov x^\top\ov v\alpha)^2]/2.


\end{equation}





В результате решения этой задачи находятся оптимальные векторы $p^*$, 


$q^*$, $\alpha^*$,


по которым вычисляются невязки несовместной системы


$$


w^*_x=(A^\top p^*+\ov v\alpha^*)_+,\quad


w^*_v=(K^\top q^*+\ov x\alpha^*)_+,\quad


w^*_3=\rho-b^\top p^*-d^\top q^*-\ov x^\top\ov v\alpha^*.


$$


С помощью этих формул в соответствии с теоремой 4 определяются нормальные


решения системы (51)


\begin{equation}


\wt x^*=w^*_x/w^*_3,\quad \wt v^*=w^*_v/w^*_3,


\end{equation}


где $w^*_3>0$ в силу (37), т.\,е. получены нормальные решения задач 


$(\p)$ и $(\C)$.





Решение задачи ЛП, таким образом, свелось к однократной безусловной


минимизации выпуклой дифференцируемой кусочно-квадратичной функции от 


$n+1$


переменной.





Ранее (см. формулы (29) в [12]) были введены необходимые и


достаточные условия оптимальности задачи ЛП, которые в


обозначениях этой статьи записываются в виде системы


\begin{equation}


A^\top p-\ov v\le0_n,\quad K^\top q-\ov x\le0_n,\quad


d^\top q+b^\top p=\ov x^\top \ov v.


\end{equation}


Найдя произвольное решение этой системы, по формулам


$$


x=\ov x-K^\top q,\quad v=\ov v-A^\top p


$$


получаем решение задач $(\p)$ и $(\C)$.





Очевидно, если положить $\alpha=-1$, $\rho=0$, то (52) переходит в 


систему


(55). Принципиальное отличие (52) от (55) состоит в том, что


система (55) совместна, а система (52) не совместна. При


минимизации невязки в системе (52) приходится проводить


безусловную минимизацию по переменным, число которых на единицу


больше, чем при минимизации невязки в (55). Зато в результате


минимизации невязки несовместной системы (52) получаются


единственные, нормальные решения задач $(\p)$ и $(\C)$, в то время как


из минимизации невязки системы (55) получаются какие-либо из


возможных решений $(\p)$ и $(\C)$.





В [12] для отыскания нормальных решений задач $(\p)$ и $(\C)$ предлагалось


использовать решение задачи


\begin{equation}


\min_{p\in R^m}\;\min_{q\in R^\nu}\;\min_{\alpha\in R^1}


[\|(A^\top p-\ov v\alpha)_+\|^2+


\|(K^\top q-\ov x\alpha)_+\|^2-


b^\top p-d^\top q+\ov x^\top\ov v\alpha].


\end{equation}


Пусть $\wh p$, $\wh q$, $\wh\alpha$ --- решения этой задачи. Из 


необходимых и


достаточных условий оптимальности задачи (56) следуют формулы


для нормального решения системы (51)


$$


\wt x^*=(A^\top\wh p-\ov v\wh\alpha)_+,\quad


\wt v^*=(K^\top\wh q-\ov x\wh\alpha)_+.


$$


Сравнивая эти выражения с (54), приходим к формулам $\wh p=p^*/w^*_3$, 


$\wh q=q^*/w^*_3$, $\wh\alpha=-\alpha^*/w^*_3$.





Таким образом, оба подхода весьма близки. В варианте (53) дополнительно


появляется произвольный положительный параметр $\rho$, несущественно 


изменяющий


метод.





\begin{thebibliography}{99}





\bibitem{1}


Еремин И.И., Мазуров В.Д., Астафьев Н.Н. {\em Несобственные задачи


линейного и выпуклого программирования}. -- М.: Наука, 1983. -- 335 с.


\bibitem{2}


Еремин И.И. {\em Противоречивые модели оптимального планирования}. -- М.:


Наука, 1988. -- 160~с.


\bibitem{3}


Еремин И.И. {\em Теория линейной оптимизации}. -- Екатеринбург: Изд-во


``Екатеринбург", 1999. -- 247 с.


\bibitem{4}


Еремин И.И. {\em Двойственность в линейной оптимизации}. -- Екатеринбург:


УрО РАН, 2001. -- 179 с.


\bibitem{5}


Гейл Д. {\em Теория линейных экономических моделей}. -- М.: Ин. лит.,


1963. -- 418 с.


\bibitem{6}


Mangasarian O.L. {\em Nonlinear programming}. -- Philadelphia:


SIAM, 1994. -- 220 p.


\bibitem{7}


Разумихин Б.С. {\em Физические модели и методы теории равновесия в


программировании и экономике}. -- М.: Наука, 1975. -- 304 с.


\bibitem{8}


Dax A. {\em The relationship between theorems of the alternative,


least norm problems, steepest descent directions and degeneracy$:$


a review} // Annals of Operations Research. -- 1993. -- V.\,46. --


\No\,1. -- P.\,11--60.


\bibitem{9}


Васильев Ф.П., Иваницкий А.Ю. {\em Линейное программирование}. -- М.:


Факториал, 1998. -- 176~с.


\bibitem{10}


Схрейвер А. {\em Теория линейного и целочисленного программирования}. --


М.: Мир, 1991. -- 360~с.


\bibitem{11}


Еремин И.И. {\em О квадратичных и полноквадратичных задачах


выпуклого программирования} // Изв. вузов.


Математика. -- 1998. -- \No\,12. -- С.\,22--28.


\bibitem{12}


Голиков А.И., Евтушенко Ю.Г. {\em Отыскание нормальных решений в задачах


линейного программирования} // Журн. вычисл. матем. и матем. физ.


-- 2000. -- Т.\,40. -- \No\,12. -- С.\,1766--1786.





\end{thebibliography}





\vskip 1cm





\post{\it Вычислительный центр}{\it Поступила}


\post{\it Российской Академии наук}{29.06.2001}





\label{end}


\end{document}








